% =============================================================
% Capitulo 8 - Desarrollo del firmware del microcontrolador
% =============================================================
\chapter{Desarrollo del firmware del microcontrolador}
\label{cap:firmware}

\section{Arquitectura del software embebido}
\label{sec:arquitectura_firmware}

El firmware del microcontrolador se desarrolla en C++ utilizando el SDK
oficial de Raspberry Pi Pico \cite{sat2026,informeMarzo2026}. El c\'odigo
fuente y los avances de desarrollo se documentan y versionan en el
repositorio del proyecto \cite{kicad_hsrl}.

% PLACEHOLDER: describir la arquitectura de software definitiva
% (maquina de estados, uso de interrupciones/DMA para el ADC,
% organizacion de modulos) una vez avanzado el desarrollo del firmware.
% Referenciar el codigo fuente relevante en el Anexo
% \ref{anexo:codigo}.

\section{Adquisici\'on y procesamiento de se\~nales}
\label{sec:adquisicion_procesamiento}

La adquisici\'on de las se\~nales $P_+$ y $P_-$ se realiza digitalizando
la salida del circuito detector de picos descrito en el
Cap\'itulo~\ref{cap:diseno_deteccion_picos} mediante el ADC integrado
del RP2350, siguiendo el criterio de tiempo de conversi\'on presentado
en la Secci\'on~\ref{sec:adc_pulsadas} \cite{informeEneFeb2026}.

% PLACEHOLDER: detallar la rutina de adquisicion implementada (canales
% utilizados, promediado de mediciones, filtrado digital si
% corresponde).

\section{Implementaci\'on del algoritmo de control}
\label{sec:algoritmo_control}

El firmware implementa la ley de control proporcional presentada en la
Secci\'on~\ref{sec:control_proporcional} (Ecuaciones~\ref{eq:error_control}
y~\ref{eq:ley_control}), calculando en cada iteraci\'on la se\~nal de
error $e(k)$ a partir de los valores de pico adquiridos y actualizando
la temperatura de consigna del seeder $T(k)$ \cite{informeMarzo2026}.

% PLACEHOLDER: incorporar el valor de Kp utilizado y el detalle de
% implementacion (frecuencia de ejecucion del lazo, saturacion de la
% salida, condicion de convergencia) una vez definidos experimentalmente.

\section{Comunicaci\'on con el seeder y con la PC}
\label{sec:comunicacion_seeder_pc}

La comunicaci\'on con el l\'aser seeder se realiza mediante el enlace
RS232 descrito en la Secci\'on~\ref{sec:comunicacion_seeder}, utilizando
los comandos especificados en el manual del fabricante (Anexo~\ref{anexo:manual_seeder}).
Dicha comunicaci\'on fue verificada de acuerdo a lo especificado en ese
manual \cite{informeAbril2026}. La comunicaci\'on con la PC se realiza
por USB \cite{sat2026}; al momento de redacci\'on de este informe, se
encuentra funcional el env\'io de datos desde la placa hacia la PC
(estado de sintonizaci\'on, par\'ametros), restando la implementaci\'on
del env\'io de comandos desde la PC hacia la placa \cite{informeAbril2026}.

\section{Modos de operaci\'on del sistema}
\label{sec:modos_operacion}

El firmware del microcontrolador contempla tres modos b\'asicos de
funcionamiento \cite{sat2026}:

\begin{enumerate}
  \item \textbf{Modo operativo}: mantiene el sistema sintonizado en una
        longitud de onda deseada, ejecutando el lazo de control descrito
        en la Secci\'on~\ref{sec:algoritmo_control}.
  \item \textbf{Modo barrido}: realiza un barrido entre dos temperaturas
        de cavidad del seeder definidas, con el fin de caracterizar la
        absorci\'on del filtro rechaza-banda de iodo.
  \item \textbf{Modo de parada}: detiene el funcionamiento del sistema.
\end{enumerate}

PLACEHOLDER: documentar el estado de implementaci\'on de cada modo de
operaci\'on y, cuando corresponda, las m\'aquinas de estado o diagramas
de transici\'on utilizados en el firmware.
